%%%%%%%%%%%%%%%%%%%%%%%%%%%%%%%%%%%%%%%%%%%%%%%%%%%%%%%%%%%%
% Preamble:
%%%%%%%%%%%%%%%%%%%%%%%%%%%%%%%%%%%%%%%%%%%%%%%%%%%%%%%%%%%%

% KOMA article class
% see also:
% http://mirror.ctan.org/macros/latex/contrib/koma-script/scrguide.pdf
\documentclass[english]{scrartcl}

\usepackage[utf8]{inputenc}
\usepackage[T1]{fontenc}
\usepackage[final]{microtype}
\usepackage{fancyhdr}
\pagestyle{fancy}

% use more space for the text (a4)
\renewcommand{\footskip}{1cm}
\renewcommand{\textheight}{690pt}
\renewcommand{\textwidth}{400pt}

% latin modern - 'The fonts, as compared to the CM family,
% contain a lot of additional characters, mainly accented
% ones.' http://www.ctan.org/tex-archive/fonts/lm/
\usepackage{lmodern}

% Linemargin
\usepackage{setspace}
\onehalfspacing

\usepackage[english]{babel}
% \usepackage[ngerman]{babel}

% e.g.: \includegraphics ...
\usepackage{graphicx}
\usepackage{subcaption}
\usepackage{float}

% e.g. \align ...
\usepackage{amsmath}
% prettier fonts
% \usepackage{mathpazo}

% nicer tables:
% \usepackage{colortbl}
\usepackage{booktabs}

% code
\usepackage{listings}
\usepackage{color}

\lstset{
	%backgroundcolor=\color[rgb]{0.95, 0.95, 0.95},
	xleftmargin=20pt,
	tabsize=2,
	rulecolor=,
	basicstyle=\footnotesize\ttfamily,
	aboveskip={1\baselineskip},
	columns=fixed,
	showstringspaces=false,
	extendedchars=true,
	breaklines=true,
	prebreak = \raisebox{0ex}[0ex][0ex]{\ensuremath{\hookleftarrow}},
	frame=l,
	numbers=left,
	numberstyle=\scriptsize\ttfamily,
	showtabs=false,
	showspaces=false,
	showstringspaces=false,
	identifierstyle=\ttfamily,
	keywordstyle=\color[rgb]{1.0,0,0},
	keywordstyle=[1]\color[rgb]{0,0,0.75},
	keywordstyle=[2]\color[rgb]{1,0.0,0.5},
	keywordstyle=[3]\color[rgb]{0.127,0.427,0.514},
	keywordstyle=[4]\color[rgb]{0.4,0.4,0.4},
	commentstyle=\color[rgb]{0.5,0.5,0.5},%\color[rgb]{0.133,0.545,0.133},
	stringstyle=\color[rgb]{0.639,0.082,0.082},
}

% special characters in biblio
\usepackage{bibgerm}

% advanced tables
\usepackage{bigdelim}
\usepackage{multirow}
\usepackage{longtable}

% rotate
\usepackage{rotating}
\def\rot#1{\begin{sideways}#1\end{sideways}}

%paragraph formatting
\usepackage{parskip}

% always last:
% e.g. \url, \href, inter-PDF-links ...
\usepackage{hyperref}

% path for graphics
\graphicspath{{./img/}}

% defs


\fancyhead[C]{\textbf{BIIGLE On The Ship How-To}}

\begin{document}
\subsection*{Access BIIGLE on the BIIGLE laptop}

\begin{enumerate}
	\item Power up the BIIGLE laptop.
	\item Log in as user BIIGLE with password \texttt{biigle}.
	\item Double click the icon named BIIGLE on the desktop. This will open BIIGLE in the Firefox browser.
\end{enumerate}

\subsection*{Access BIIGLE on another computer}

To access BIIGLE on another computer you have to set up a network connection to the BIIGLE laptop. This requires the following one-time setup:

\begin{enumerate}
	\item Power up the BIIGLE laptop.
	\item Log in as user BIIGLE with password \texttt{biigle}.
	\item Connect the BIIGLE laptop to the network. This can be done via LAN cable or WiFi. \begin{itemize}
		\item LAN: Plug in the network cable.
		\item WiFi: Open the networking menu in the upper right title bar (\raisebox{-2mm}{\includegraphics[scale=1]{connect_established.pdf}} or \raisebox{-2mm}{\includegraphics[scale=1]{connect_no.pdf}}), click on the WiFi network to connect to and enter the network credentials in the popup window.
	\end{itemize}
	\item Find out the BIIGLE network address. Once the BIIGLE laptop is connected to the network, the address is displayed on the desktop. Example address: \texttt{http://192.168.100.1/}.
\end{enumerate}

To access BIIGLE on another computer, leave the BIIGLE laptop up and running. The other computer needs to be connected to the same network than the BIIGLE laptop. On the other computer, open a browser such as Firefox or Chrome and enter the BIIGLE network address. This will bring up the BIIGLE login screen.

\subsection*{Create a new BIIGLE user}

New BIIGLE users can only be created by BIIGLE administrators.

\begin{enumerate}
	\item Log in to BIIGLE as an andministrator.
	\item Open the menu (\raisebox{-0.5mm}{\includegraphics[height=0.4cm]{bars.pdf}}) and click on \textbf{Admin area}.
	\item Go to the \textbf{Users} page and click on the \textbf{New user} button at the top.
	\item Enter the user information and click \textbf{Create}.
\end{enumerate}

\subsection*{Create a new volume with images on the BIIGLE laptop}

The BIIGLE laptop has 1 TB of storage capacity for images. Images can be transferred from an external USB drive to the BIIGLE laptop as follows:

\begin{enumerate}
	\item Log in as user BIIGLE with password \texttt{biigle}.
	\item Open the \texttt{images-local} directory on the desktop.
	\item Plug in the USB drive. A second window will appear that shows the contents of the USB drive.
	\item Copy the image directories from the USB drive to the \texttt{images-local} directory.
\end{enumerate}

Now create a new volume:

\begin{enumerate}
	\item In BIIGLE, click on the ``browse'' button of the volume URL input.
	\item Choose the \texttt{local} storage disk from the select menu. This will open the directory browser of all folders in the \texttt{images-local} directory.
	\item Choose one for the new volume with a click on the \raisebox{-0.5mm}{\includegraphics[height=0.4cm]{check.pdf}} button.
\end{enumerate}

Do not use the BIIGLE laptop as a backup solution for your images!

\subsection*{Create a new volume with images on an external USB drive}

BIIGLE can also read images directly from a USB drive connected to the BIIGLE laptop. Whenever you connect a new USB drive, you have to log in once as user BIIGLE with password \texttt{biigle}, so the drive is detected.

\begin{enumerate}
	\item In BIIGLE, click on the ``browse'' button of the volume URL input.
	\item Choose the \texttt{media} storage disk from the select menu. This will open the directory browser of all USB drives connected to the BIIGLE laptop.
	\item Select one drive and choose a directory for the new volume with a click on the \raisebox{-0.5mm}{\includegraphics[height=0.4cm]{check.pdf}} button.
\end{enumerate}

Make sure to have the USB drive connected whenever you want to work with the images in BIIGLE (not only when you create a volume).

\end{document}


